\documentclass[12pt,a4paper]{article}
\usepackage{fontspec}
\setmainfont{DejaVu Serif}
\usepackage[russian]{babel}
\usepackage{geometry}
\usepackage{amsmath,amssymb,mathtools}
\usepackage{array,booktabs,longtable}
\usepackage{xcolor,hyperref,tikz}
\geometry{left=22mm,right=20mm,top=22mm,bottom=23mm}
\hypersetup{colorlinks=true,linkcolor=blue!55!black}
\setlength{\parindent}{0pt}
\setlength{\parskip}{6pt}
\renewcommand{\arraystretch}{1.3}
\newcommand{\parthead}[1]{\par\medskip\textbf{#1}\par\smallskip}

\begin{document}
\begin{titlepage}
\thispagestyle{empty}
\vspace*{0.28\textheight}
\begin{center}
{\LARGE\bfseries Ревью домашней работы\par}
\vspace{16mm}
{\Large Ксения Клыкова\par}
\vspace{9mm}
{\large 22 октября 2026 года\par}
\vspace{12mm}
{\large Теория чисел: составность и делимость\par}
\end{center}
\vfill
\begin{center}
Лёвин Артём Александрович
\end{center}
\end{titlepage}

\section*{Сводная таблица ревью}
В документ включены только задания с ошибками, неполными доказательствами или без решения. Правильно выполненные задания опущены.

\begin{longtable}{@{}p{1.1cm}p{3.1cm}p{10.6cm}@{}}
\toprule
\textbf{№} & \textbf{Тема} & \textbf{Суть замечания}\\
\midrule
\endfirsthead
\toprule
\textbf{№} & \textbf{Тема} & \textbf{Суть замечания}\\
\midrule
\endhead
\hyperref[task:5]{5} & Составные числа & Алгебраическое разложение получено верно. Проверен один частный случай вместо произвольного натурального числа; при подстановке допущена ошибка в арифметике.\\ \midrule
\hyperref[task:6]{6} & Натуральные множители & Найдены верные корни квадратного трёхчлена, но разложение с дробными множителями ещё не доказывает составность.\\ \midrule
\hyperref[task:9]{9} & Делимость на шесть & После вынесения общего множителя ошибочно применена формула разности квадратов; доказательство для целых чисел не завершено.\\ \midrule
\hyperref[task:10]{10} & Разность степеней & Решение задания на фотографии отсутствует.\\
\bottomrule
\end{longtable}

\parthead{Главный критерий}
Чтобы доказать составность натурального числа, надо предъявить его разложение на два натуральных множителя, каждый больше единицы. Чтобы доказать утверждение «для каждого $n$», нужно провести рассуждение с произвольным допустимым $n$. Проверка одного или нескольких чисел этого не заменяет.

\clearpage
\phantomsection\label{task:5}
\section*{Задание № 5. Доказательство составности}
\parthead{Текст задания}
Докажите, что $n^2+7n+10$ --- составное число при каждом $n\in\mathbb N$.

\parthead{Ваше решение}
Вы правильно разложили трёхчлен:
\[
n^2+7n+10=(n+5)(n+2).
\]
Затем рассмотрели $n=5$ и записали результат $5\cdot7=35$.

\parthead{Ошибка}
\textcolor{red}{Первое замечание:} ошибка при подстановке. Если $n=5$, то $n+5=10$; следовательно,
\[
(5+5)(5+2)=10\cdot7=70,
\]
а не $35$.

\textcolor{red}{Второе замечание:} проверка единственного значения $n=5$ подтверждает составность одного конкретного числа. Утверждение сформулировано для \emph{каждого натурального} $n$.

\parthead{Где заканчивается верное рассуждение}
Равенство $(n+5)(n+2)$ найдено верно. Необходимо продолжить его доказательством того, что оба множителя натуральные и больше единицы \emph{при любом} $n\in\mathbb N$.

\parthead{\textcolor{blue!60!black}{Ключевая идея}}
Завершить уже полученное разложение оценкой каждого множителя.

\parthead{Полное правильное доказательство}
Сгруппируем слагаемые:
\[
\begin{aligned}
n^2+7n+10
&=n^2+5n+2n+10\\
&=n(n+5)+2(n+5)\\
&=(n+5)(n+2).
\end{aligned}
\]
Пусть $n$ --- произвольное натуральное число. Тогда $n\ge1$, а значит,
\[
n+5\ge6>1,\qquad n+2\ge3>1.
\]
Числа $n+5$ и $n+2$ натуральные; каждый множитель больше единицы. Поэтому произведение $(n+5)(n+2)$ по определению является составным числом. Рассуждение применимо ко всем $n\in\mathbb N$, что и требовалось доказать.

\parthead{Проверка}
Для $n=5$ получаем:
\[
5^2+7\cdot5+10=25+35+10=70=10\cdot7.
\]
Это правильная проверка частного случая, дополняющая общее доказательство.

\parthead{\textcolor{teal!60!black}{Итог}}
Идея разложения на множители верная. Исправляем числовой пример и добавляем универсальное обоснование составности.

\parthead{Задача на закрепление}
Докажите, что $n^2+9n+20$ --- составное число при каждом $n\in\mathbb N$.

\clearpage
\phantomsection\label{task:6}
\section*{Задание № 6. Доказательство составности}
\parthead{Текст задания}
Докажите, что $6n^2+11n+3$ --- составное число при каждом $n\in\mathbb N$.

\parthead{Ваше решение}
Вы верно вычислили дискриминант:
\[
D=11^2-4\cdot6\cdot3=121-72=49
\]
и нашли корни $-\frac32$ и $-\frac13$. Из них действительно следует
\[
6n^2+11n+3=
6\left(n+\frac32\right)\left(n+\frac13\right).
\]
Далее Вы попытались обосновать составность положительностью множителей.

\parthead{Ошибка}
\textcolor{red}{Переход к выводу о составности недостаточно обоснован.}
Составное число по определению должно быть произведением двух \emph{натуральных} множителей, каждый из которых больше единицы. Выражения $n+\frac32$ и $n+\frac13$ при натуральном $n$ дробные. Положительность множителей ничего не говорит об их натуральности. Даже простое число $7$ можно записать в виде произведения двух положительных дробных или целых чисел больше единицы:
\[
7=\frac72\cdot2.
\]
Поэтому из дробного разложения нельзя сразу делать вывод о составности.

\parthead{Исправление от последнего правильного шага}
\textcolor{blue!60!black}{Ключевая идея:} перенести знаменатели дробей в коэффициенты линейных множителей:
\[
\begin{aligned}
6\left(n+\frac32\right)\left(n+\frac13\right)
&=6\cdot\frac{2n+3}{2}\cdot\frac{3n+1}{3}\\
&=(2n+3)(3n+1).
\end{aligned}
\]
Получено разложение с целыми коэффициентами.

\parthead{Полное правильное доказательство}
Его можно получить и без дискриминанта, группировкой:
\[
\begin{aligned}
6n^2+11n+3
&=6n^2+9n+2n+3\\
&=3n(2n+3)+(2n+3)\\
&=(3n+1)(2n+3).
\end{aligned}
\]
Если $n\in\mathbb N$, то
\[
3n+1\ge4>1,\qquad 2n+3\ge5>1.
\]
Оба множителя натуральные при каждом натуральном $n$. Поэтому их произведение составно при каждом $n\in\mathbb N$.

\parthead{Проверка}
При $n=1$:
\[
6+11+3=20=(3\cdot1+1)(2\cdot1+3)=4\cdot5.
\]
Раскрывая скобки $(3n+1)(2n+3)$, возвращаем исходный трёхчлен.

\parthead{\textcolor{teal!60!black}{Итог}}
Дискриминант и найденные корни правильные. Недостающий ключевой шаг --- переход к натуральным множителям, каждый из которых больше единицы.

\parthead{Задача на закрепление}
Докажите, что $6n^2+13n+6$ --- составное число при каждом $n\in\mathbb N$.

\clearpage
\phantomsection\label{task:9}
\section*{Задание № 9. Делимость на шесть}
\parthead{Текст задания}
Докажите, что $n^5-n$ делится на $6$ для каждого $n\in\mathbb Z$.

\parthead{Ваше решение}
Вы начали с вынесения общего множителя $n$ и затем пытались применить формулу разности квадратов. В последней читаемой строке записаны дробные показатели степени $\frac52$; полученное произведение не соответствует исходному выражению. Завершённого доказательства делимости на $6$ нет.

\parthead{Ошибка}
\textcolor{red}{Ошибка в выборе степеней при разложении.}
После вынесения $n$ остаётся $n^4-1$:
\[
n^5-n=n(n^4-1).
\]
Поскольку $n^4=(n^2)^2$, правильная формула даёт
\[
n^4-1=(n^2-1)(n^2+1).
\]
Степень $\frac52$ здесь не возникает. Кроме того, дробные степени вида $n^{5/2}$ не определены как действительные числа для отрицательных $n$, а задача требует рассмотреть \emph{все целые} $n$.

\parthead{Исправление от последнего правильного шага}
\textcolor{blue!60!black}{Ключевая идея:} выделить произведение трёх последовательных целых чисел:
\[
\begin{aligned}
n^5-n
&=n(n^4-1)\\
&=n(n^2-1)(n^2+1)\\
&=(n-1)n(n+1)(n^2+1).
\end{aligned}
\]

\parthead{Полное правильное доказательство}
Для любого целого $n$ числа $n-1$, $n$, $n+1$ --- три последовательных целых числа.

\begin{enumerate}
\item Среди любых двух последовательных целых чисел одно чётное. Поэтому произведение $(n-1)n(n+1)$ делится на $2$.
\item Среди любых трёх последовательных целых чисел одно делится на $3$. Поэтому то же произведение делится на $3$.
\item Поскольку $2$ и $3$ взаимно просты, число, делящееся одновременно на $2$ и $3$, делится на $6$.
\end{enumerate}
Следовательно, существует целое число $k$ такое, что
\[
(n-1)n(n+1)=6k.
\]
Отсюда
\[
n^5-n=6k(n^2+1).
\]
Число $k(n^2+1)$ целое. Значит,
\[
\boxed{6\mid(n^5-n)\quad\text{при каждом }n\in\mathbb Z.}
\]

\parthead{Проверка}
При $n=-2$ получаем $(-2)^5-(-2)=-30$, а при $n=0$ получаем $0$. Оба числа делятся на $6$, и общее доказательство охватывает эти случаи.

\parthead{\textcolor{teal!60!black}{Итог}}
Вынесение $n$ было правильным началом. После него следует разложить четвёртую степень и воспользоваться свойством трёх последовательных целых чисел.

\parthead{Задача на закрепление}
Докажите, что $n^7-n$ делится на $6$ при каждом $n\in\mathbb Z$.

\clearpage
\phantomsection\label{task:10}
\section*{Задание № 10. Делимость на пятнадцать}
\parthead{Текст задания}
Докажите, что $2^{4n}-1$ делится на $15$ при каждом $n\in\mathbb N$.

\parthead{Ваше решение}
На представленной фотографии решение задания № 10 отсутствует. Неверных преобразований нет, поскольку задача пока не выполнена. Разберём её с первого шага.

\parthead{Что необходимо доказать}
Нужно выделить множитель $15$ при произвольном натуральном $n$. Проверка отдельных значений не заменит общего доказательства.

\parthead{\textcolor{blue!60!black}{Ключевая идея}}
Заметить, что $2^4=16$ и $16-1=15$, после чего применить школьную формулу разности степеней:
\[
a^n-b^n=(a-b)(a^{n-1}+a^{n-2}b+\dots+b^{n-1}).
\]

\parthead{Полное правильное доказательство}
Для любого натурального $n\ge1$ имеем
\[
2^{4n}-1=(2^4)^n-1=16^n-1.
\]
По формуле разности степеней:
\[
\begin{aligned}
16^n-1
&=(16-1)(16^{n-1}+16^{n-2}+\dots+16+1)\\
&=15(16^{n-1}+16^{n-2}+\dots+16+1).
\end{aligned}
\]
Выражение в скобках --- целое число, так как это сумма целых степеней числа $16$. Поэтому всё выражение делится на $15$. Доказательство верно для любого натурального $n$.

\parthead{Проверка}
При $n=1$:
\[
2^4-1=15.
\]
При $n=2$:
\[
2^8-1=255=15\cdot17.
\]
Эти примеры согласуются с уже выполненным общим доказательством.

\parthead{\textcolor{teal!60!black}{Итог}}
Достаточно перейти к степени числа $16$ и вынести множитель $16-1=15$.

\parthead{Задача на закрепление}
Докажите, что $3^{4n}-1$ делится на $16$ при каждом $n\in\mathbb N$.

\vfill
\begin{center}
\small Лёвин Артём Александрович эксклюзивно для Ксении Клыковой
\end{center}
\end{document}
